%==============================================================================
% Sjabloon poster bachproef
%==============================================================================
% Gebaseerd op document class `a0poster' door Gerlinde Kettl en Matthias Weiser
% Aangepast voor gebruik aan HOGENT door Jens Buysse en Bert Van Vreckem

\documentclass[a0,portrait]{hogent-poster}

% Info over de opleiding
\course{Bachelorproef}
\studyprogramme{toegepaste informatica}
\academicyear{2025-2026}
\institution{Hogeschool Gent, Valentin Vaerwyckweg 1, 9000 Gent}

% Info over de bachelorproef
\title{Analyse van risico's bij IBM DB2 transaction logreplicatie van kritieke bancaire mainframe-systemen naar geodistribueerde disaster recovery\--omge\-vingen, met focus op RPO en RTO}
% \subtitle{Ondertitel (eventueel)}
\author{Robin Ledoux}
\email{robin.ledoux@student.hogent.be}
\supervisor{Leendert Blondeel}
\cosupervisor{Jan Cannaerts (Rabobank)}

% Indien ingevuld, wordt deze informatie toegevoegd aan het einde van de
% abstract. Zet in commentaar als je dit niet wilt.
\specialisation{mainframe}
\keywords{mainframe, Db2, logreplicatie, geodistribueerde systemen, RPO, RTO, disaster recovery}
\projectrepo{https://github.com/RobinLedoux/latex-hogent-bachelorproef}

\begin{document}

\maketitle

\begin{abstract}
Financiële instellingen maken gebruik van synchrone data-mirroring tussen twee geografisch nabijgelegen mainframe-datacenters om hoge beschikbaarheid en transactionele consistentie te waarborgen.
Bij gelijktijdige uitval van beide sites, bijvoorbeeld door grootschalige regionale incidenten, volstaat deze opstelling echter niet om de continuïteit van kritieke diensten en de integriteit van transactiedata te garanderen.
Om dit risico te minimaliseren wordt in dit onderzoek onderzocht hoe een derde, geografisch gescheiden site kan worden ingezet binnen een IBM Db2 voor z/OS-omgeving.
Het onderzoek wordt uitgevoerd aan de hand van een technische literatuurstudie, een vergelijkende architecturale analyse en een risicoanalyse.
Hierbij worden synchrone en asynchrone replicatietechnieken voor datareplicatie zoals IBM GDPS Metro Mirror en IBM GDPS Global Mirror onderzocht, samen met IBM Q Replication voor logreplicatie binnen Db2 voor z/OS.
De analyse focust op transactionele consistentie, herstelmechanismen en de impact op Recovery Point Objective (RPO) en Recovery Time Objective (RTO).
De resultaten tonen aan dat het toevoegen van een derde asynchrone site onvermijdelijk bijkomende risico's introduceert, met name op het vlak van langdurige transacties, verhoogde RPO en complexere herstelprocedures.
Tegelijkertijd blijkt dat deze risico's beheersbaar zijn door gerichte technische maatregelen en goed gedefinieerde herstelprocedures.
Dit onderzoek biedt een gestructureerd basis dat financiële instellingen ondersteunt bij het maken van geïnformeerde keuzes rond geodistribueerde architecturen in Db2 voor z/OS-omgevingen bij herstelscenario's.
\end{abstract}

\begin{multicols}{2} % This is how many columns your poster will be broken into, a portrait poster is generally split into 2 columns

\section{Introductie}

Financiële instellingen gebruiken synchrone data-mirroring voor twee geografisch nabijgelegen datacenters om de beschikbaarheid van hun mainframe-omgevingen te waarborgen. 
Europese voorschriften stellen strikte eisen aan de operationele veerkracht en continuïteit van financiële instellingen.
De Digital Operational Resilience Act verplicht financiële instellingen om een uitgebreid ICT-risico\-beheer\-kader te implementeren dat de beschikbaarheid en integriteit van kritieke systemen waarborgt, inclusief plannen voor redundantie of herstel wanneer er zich verstoringen voordoen. 

Het probleem ontstaat wanneer twee gesynchroniseerde datacenters tegelijkertijd uitvallen, waardoor de continuïteit van de diensten en de integriteit van transacties in gevaar komt. 
Dit scenario is relevant bij grootschalige regionale incidenten zoals stroomuitval, natuurrampen of gecoördineerde cyberaanvallen.
Het toevoegen van een aanvullende derde site zorgt voor uitdagingen omdat door de afstand geen minimale vertragingstijd kan worden gegarandeerd.
Om dit risico te minimaliseren is het noodzakelijk om een extra, geografisch gescheiden datacentrum toe te voegen, waarop alle transacties asynchroon kunnen worden gerepliceerd.
Een derde aanvullende site betekent echter niet dat er geen invloed is op de vertragingstijd.
Hierdoor is onderzoek naar verschillende oplossingen met betrekking tot een derde aanvullende site bij een oplossing van reeds twee gesynchroniseerde datacenters nuttig om een overzicht te krijgen van de invloed op Recovery Point Objective (RPO) en Recovery Time Objective (RTO).

Deze bachelorproef beoogt de volgende deelvragen te beantwoorden.

\begin{itemize}
    \item Welke vormen van datareplicatie en Db2-logreplicatie bestaan binnen Db2 voor z/OS en welke garanties bieden deze inzake transactionele consistentie?
    \item Welke technische beperkingen (vertragingstijd, afstand, commit-gedrag) verhinderen het gebruik van synchrone replicatie naar een derde geografisch afgelegen site?
    \item Welke faalscenario's kunnen optreden in een Db2-architectuur met meerdere sites en asynchrone replicatie, en welke Db2-componenten zijn daarbij betrokken?
    \item Welke impact hebben deze faalscenario's op dataconsistentie, RPO en RTO\@?
\end{itemize}

\section{Risicomatrix}

\small
\begin{tabularx}{\columnwidth}{|
  >{\raggedright\arraybackslash}p{6.5cm} | % Risico (korter)
  p{3cm} | % Waarschijnlijkheid (korter)
  p{3cm} | % Impact (korter)
  p{2cm} | % Score (korter)
  >{\raggedright\arraybackslash}X | % Toelichting (flexibel)
}
\hline
\textbf{Risico} & \textbf{Waar\-schijn\-lijk\-heid (1--3)} & \textbf{Impact (1--3)} & \textbf{Score} & \textbf{Toelichting} \\
\hline
Replicatie\-vertraging & 3 & 2 & 6 & Inherent aan asynchrone replicatie; frequentie stijgt bij hoge workload of beperkte bandbreedte. \\
\hline
Long-running LUW's & 2 & 3 & 6 & Minder frequent maar met hoge impact op dataconsistentie en RPO. \\
\hline
Inconsistenties in volgorde & 2 & 3 & 6 & Kan leiden tot foutieve saldi of dubbele transacties bij financiële instellingen. \\
\hline
Geregionaliseerde uitval (twee synchrone sites) & 1 & 3 & 3 & Beperkte kans, maar catastrofale gevolgen als beide sites in hetzelfde rampgebied liggen. \\
\hline
Netwerk\-onderbreking replicatie & 2 & 2 & 4 & Redundante verbindingen kunnen dit risico beperken. \\
\hline
Gaten in logbestanden & 1 & 3 & 3 & Zeldzaam maar met directe impact op RPO en complexiteit van herstel. \\
\hline
Queue-overbelasting (MQ) & 2 & 2 & 4 & Verhoogde workload kan MQ-queues overbelasten en replicatievertraging vergroten. \\
\hline
Failover naar asynchrone site & 2 & 2 & 4 & Vereist complete replicatie van consistency groups voordat site C operationeel kan zijn. \\
\hline
\end{tabularx}

\section{Conclusies}

Uit de analyse blijkt dat een opstelling met uitsluitend twee synchrone datacenters, hoewel deze een RPO van nul garandeert, onvoldoende bescherming biedt tegen geregionaliseerde uitval.
De afstandsbeperkingen inherent aan synchrone replicatie maken het onmogelijk om volledige geografische spreiding te realiseren zonder significante impact op de prestaties.
Asynchrone replicatietechnologieën, zoals GDPS Global Mirror, bieden daarentegen wel de mogelijkheid tot geografische spreiding over grote afstanden, maar introduceren onvermijdelijk risico's op dataverlies en inconsistentie. 
Deze risico's manifesteren zich voornamelijk onder de vorm van replicatievertraging, incomplete Logical Units of Work (LUW's) en inconsistenties in de volgorde van transacties. 
Deze fenomenen hebben een directe negatieve impact op zowel RPO als RTO en vereisen bijkomende herstelmechanismen.
De hybride architectuur, waarbij twee synchrone sites worden gecombineerd met een derde asynchrone site, vormt op basis van dit onderzoek de meest robuuste oplossing. 
De keerzijde is echter een verhoogde complexiteit en kost, omdat er een extra site moet worden voorzien.
De bijdrage van dit onderzoek ligt in het systematisch in kaart brengen van de interactie tussen replicatiestrategieën, faalscenario's en hun impact op RPO en RTO binnen Db2 voor z/OS-omgevingen.


\section{Toekomstig onderzoek}

Verder onderzoek kan zich richten op experimentele validatie van de beschreven architecturen, bijvoorbeeld via simulaties of testomgevingen. 
Daarnaast is er ruimte voor onderzoek naar geautomatiseerde failovermechanismen en optimalisatietechnieken die de impact van replicatievertraging en Controlled Log Catch-Up (CLC) op RTO verder kunnen beperken. 
Ook de rol van nieuwe technologieën, zoals AI-gedreven monitoring en voorspellende faal analyse, vormt een interessant perspectief binnen dit domein.
Samenvattend kan gesteld worden dat er geen universele oplossing bestaat die simultaan een RPO van nul, minimale RTO en volledige geografische spreiding garandeert. 
De keuze voor een replicatiestrategie blijft steeds een afweging tussen consistentie, performantie en operationele complexiteit. 
De hybride architectuur biedt hierbij de beste balans voor kritieke systemen binnen financiële omgevingen, mits de bijkomende risico's actief worden geminimaliseerd door de besproken mitigatiestrategieën.

\end{multicols}
\end{document}